\documentclass[11pt,letterpaper]{article}
\PassOptionsToPackage{hyphens}{url}
\usepackage[margin=.83in,top=.8in,bottom=.8in,headheight=15pt]{geometry}
\usepackage[T1]{fontenc}
\usepackage{lmodern,amsmath,amssymb,mathtools,bm,booktabs,array,graphicx,tikz}
\usetikzlibrary{arrows.meta,positioning,calc,decorations.pathreplacing,patterns}
\definecolor{ink}{HTML}{17344D}\definecolor{accent}{HTML}{087E8B}\definecolor{pale}{HTML}{EDF5F5}\definecolor{gold}{HTML}{AD6B19}
\providecommand{\dd}{\,\mathrm{d}}\providecommand{\R}{\mathbb{R}}
\providecommand{\E}{\mathbb{E}}\providecommand{\Var}{\operatorname{Var}}
\providecommand{\Cov}{\operatorname{Cov}}\providecommand{\diag}{\operatorname{diag}}
\providecommand{\norm}[1]{\left\lVert#1\right\rVert}
\providecommand{\abs}[1]{\left|#1\right|}
\usepackage{microtype,hyperref,fancyhdr,enumitem,amsthm,sectsty,needspace}
\allsectionsfont{\raggedright}
\hypersetup{colorlinks=true,linkcolor=ink,urlcolor=accent,pdfauthor={Zan Ahmad},pdftitle={Resistance networks: constitutive laws, pressure, and conservation / Review}}
\setlength{\parindent}{0pt}\setlength{\parskip}{6pt}\linespread{1.07}
\setlist{itemsep=3pt,topsep=4pt,leftmargin=*}
\emergencystretch=2em\clubpenalty=10000\widowpenalty=10000
\pagestyle{fancy}\fancyhf{}\fancyhead[L]{\small Mathematical Modeling}\fancyhead[R]{\small Lecture 2}
\fancyfoot[L]{\footnotesize Zan Ahmad}\fancyfoot[R]{\footnotesize\thepage}\renewcommand{\headrulewidth}{.3pt}
\newtheorem{proposition}{Proposition}\newtheorem{theorem}{Theorem}\newtheorem{lemma}{Lemma}
\theoremstyle{definition}\newtheorem{definition}{Definition}\newtheorem{example}{Example}
\numberwithin{equation}{section}
\begin{document}
\section*{Resistance networks: constitutive laws, pressure, and conservation --- review sheet}
Use this after attempting the worked examples and practice. Explain each claim in words before calculating.

\section*{Key ideas}
\begin{itemize}
\item Pressure difference divided by positive resistance gives signed flow.
\item A reference arrow does not enforce one-way motion.
\item Steady internal nodes have zero net accumulation.
\item Internal transfers cancel in a whole-network account.
\item Common pressure shifts preserve flow in this model.
\item Uniform resistance scaling inversely scales all flows at fixed endpoint pressures.
\end{itemize}
\section*{Common errors and repairs}
\textbf{Using one pressure instead of a difference.} Subtract the downstream reference pressure using the declared arrow.

\textbf{Clipping a negative resistor flow.} Interpret its direction; a valve requires a separate rule.

\textbf{Checking only the total boundary flow.} Also check every internal node balance.

\section*{Retrieve without your notes}
\begin{itemize}
\item What does a negative A-to-B reference flow mean?
\item What are the resistance units when pressure is mmHg and flow is L/min?
\item What does this steady model omit?
\end{itemize}
\section*{Readiness check}
\begin{itemize}
\item Use a pressure difference and positive resistance to calculate a signed flow.
\item Check that each internal node has equal total inflow and outflow.
\item Interpret a negative answer relative to a declared arrow.
\item Predict the effect of shifting the pressure reference or scaling all resistances.
\item State which steady-flow assumptions limit the model.
\end{itemize}
\end{document}
